\chapter{2007年辽宁卷（理）}

\section{单选题}

\begin{problem}
设集合 \(U = \{1, 2, 3, 4, 5\}\)，\(A = \{1, 3\}\)，\(B = \{2, 3, 4\}\)，则 \((\complement_U A) \cap (\complement_U B) =\)（\quad）

\choices
    {\(\{1\}\)}
    {\(\{5\}\)}
    {\(\{2,4\}\)}
    {\(\{1, 2, 4, 5\}\)}
\end{problem}

\begin{problem}
若函数 \( y = f(x) \) 的反函数图象过点 \((1,5)\)，则函数 \( y = f(x) \) 的图象必过点（\quad）

\choices
    {\((5,1)\)}
    {\((1,5)\)}
    {\((1,1)\)}
    {\((5,5)\)}
\end{problem}

\begin{problem}
双曲线 \(\frac{x^{2}}{16} - \frac{y^{2}}{9} = 1\) 的焦点坐标为（\quad）

\choices
    {\((-\sqrt{7},0),(\sqrt{7},0)\)}
    {\((0,-\sqrt{7}),(0,\sqrt{7})\)}
    {\(( -5,0),(5,0)\)}
    {\((0,-5),(0,5)\)}
\end{problem}

\begin{problem}
若向量 \(a\) 与 \(b\) 不共线，\(a \cdot b \neq 0\)，且 \(c = a - \left(\frac{a \cdot a}{a \cdot b}\right)b\)，则向量 \(a\) 与 \(c\) 的夹角为（\quad）

\choices
    {\(0\)}
    {\(\frac{\pi}{6}\)}
    {\(\frac{\pi}{3}\)}
    {\(\frac{\pi}{2}\)}
\end{problem}

\begin{problem}
设等差数列 \(\{a_{n}\}\) 的前 \(n\) 项和为 \(S_{n}\)，若 \(S_{3} = 9, S_{6} = 36\)，则 \(a_{7} + a_{8} + a_{9} =\)（\quad）

\choices
    {\(63\)}
    {\(45\)}
    {\(36\)}
    {\(27\)}
\end{problem}

\begin{problem}
若 \(m, n\) 是两条不同的直线，\(\alpha, \beta, \gamma\) 是三个不同的平面，则下列命题中的真命题是（\quad）

\choices
    {若 \(m \subset \beta, \alpha \perp \beta\)，则 \(m \perp \alpha\)}
    {若 \(m \perp \beta, m \parallel \alpha\)，则 \(\alpha \perp \beta\)}
    {若 \(\alpha \perp \gamma, \alpha \perp \beta\)，则 \(\beta \perp \gamma\)}
    {若 \(\alpha \cap \gamma = m, \beta \cap \gamma = n, m \parallel n\)，则 \(\alpha \parallel \beta\)}
\end{problem}

\begin{problem}
若函数 \(y = f(x)\) 的图象按向量 \(\mathbf{a}\) 平移后，得到函数 \(y = f(x - 1) - 2\) 的图象，则向量 \(\mathbf{a} =\)（\quad）

\choices
    {\((1, -2)\)}
    {\((1, 2)\)}
    {\((-1, -2)\)}
    {\((-1, 2)\)}
\end{problem}

\begin{problem}
已知变量 \(x\)、\(y\) 满足约束条件 \(\left\{ \begin{array}{l} x - y + 2 \leqslant 0, \\ x \geqslant 1, \\ x + y - 7 \leqslant 0, \end{array} \right.\) 则 \(\frac{y}{x}\) 的取值范围是（\quad）

\choices
    {\(\left[\frac{9}{5}, 6\right]\)}
    {\(\left(-\infty, \frac{9}{5}\right] \cup [6, +\infty)\)}
    {\((-\infty, 3] \cup [6, +\infty)\)}
    {\([3, 6]\)}
\end{problem}

\begin{problem}
一个坛子里有编号为 1, 2，…，12 的 12 个大小相同的球，其中 1 到 6 号球是红球，其余的是黑球。若从中任取两个球，则取到的都是红球，且至少有 1 个球的号码是偶数的概率为（\quad）

\choices
    {\(\frac{1}{22}\)}
    {\(\frac{1}{11}\)}
    {\(\frac{3}{22}\)}
    {\(\frac{2}{11}\)}
\end{problem}

\begin{problem}
设 \(p\)、\(q\) 是两个命题，\(p: \log_3(|x| - 3) > 0\)，\(q: x^2 - \frac{5}{6}x + \frac{1}{6} > 0\)，则 \(p\) 是 \(q\) 的（\quad）

\choices
    {充分而不必要条件}
    {必要而不充分条件}
    {充分必要条件}
    {既不充分也不必要条件}
\end{problem}

\begin{problem}
设 \(P\) 为双曲线 \(x^2 - \frac{y^2}{12} = 1\) 上的一点，\(F_1, F_2\) 是该双曲线的两个焦点。若 \(|PF_1| : |PF_2| = 3 : 2\)，则 \(\triangle PF_1F_2\) 的面积为（\quad）

\choices
    {\(6\sqrt{3}\)}
    {\(12\)}
    {\(12\sqrt{3}\)}
    {\(24\)}
\end{problem}

\begin{problem}
已知 \(f(x)\) 与 \(g(x)\) 是定义在 \(\R\) 上的连续函数，如果 \(f(x)\) 与 \(g(x)\) 仅当 \(x = 0\) 时的函数值为0，且 \(f(x) \geqslant g(x)\)，那么下列情形不可能出现的是（\quad）

\choices
    {0 是 \(f(x)\) 的极大值，也是 \(g(x)\) 的极大值}
    {0 是 \(f(x)\) 的极小值，也是 \(g(x)\) 的极小值}
    {0 是 \(f(x)\) 的极大值，但不是 \(g(x)\) 的极值}
    {0 是 \(f(x)\) 的极小值，但不是 \(g(x)\) 的极值}
\end{problem}

\section{填空题}

\begin{problem}
已知函数 \(y = f(x)\) 为奇函数，若 \(f(3) - f(2) = 1\)，则 \(f(-2) - f(-3) = \fillinblank{}\).
\end{problem}

\begin{problem}
设 \(\left(\sqrt{x} + \frac{1}{\sqrt[5]{x}}\right)^{8}\) 展开式中含有 \(x\) 的整数次幂的项的系数之和为 \fillinblank{}。（用数字作答）
\end{problem}

\begin{problem}
若一个底面边长为 \(\frac{\sqrt{6}}{2}\)，侧棱长为 \(\sqrt{6}\) 的正六棱柱的所有顶点都在一个球的面上，则此球的体积为 \fillinblank{}.
\end{problem}

\begin{problem}
设椭圆 \(\frac{x^2}{25} + \frac{y^2}{16} = 1\) 上一点 \(P\) 到左准线的距离为 \(10\)，\(F\) 是该椭圆的左焦点。若点 \(M\) 满足 \(\overrightarrow{OM} = \frac{1}{2}\left(\overrightarrow{OP} + \overrightarrow{OF}\right)\)，则 \(\left|\overrightarrow{OM}\right| = \fillinblank{}\).
\end{problem}

\section{解答题}

\begin{problem}
某公司在过去几年内使用某种型号的灯管 1000 支，该公司对这些灯管的使用寿命（单位：小时）进行了统计，统计结果如下表所示：

\begin{center}\small
\begin{tabular}{|c|c|c|c|c|}
\hline
分组 & \([500,900)\) & \([900,1100)\) & \([1100,1300)\) & \([1300,1500)\) \\
\hline
频数 & 48 & 121 & 208 & 223 \\
\hline
频率 &  &  &  &  \\
\hline
\end{tabular}

\vspace{0.5em}

\begin{tabular}{|c|c|c|c|}
\hline
分组 & \([1500,1700)\) & \([1700,1900)\) & \([1900,+\infty)\) \\
\hline
频数 & 193 & 165 & 42 \\
\hline
频率 &  &  &  \\
\hline
\end{tabular}
\end{center}

(1) 将各组的频率填入表中;

(2) 根据上述统计结果，计算灯管使用寿命不足 1500 小时的频率;

(3) 该公司某办公室新安装了这种型号的灯管 3 支，若将上述频率作为概率，试求至少有 2 支灯管的寿命不足 1500 小时的概率.
\end{problem}

\begin{problem}
如图，在直三棱柱 \(ABC - A_{1}B_{1}C_{1}\) 中，\(\angle ACB = 90^{\circ}\)，\(AC = BC = a\)，\(D\)，\(E\) 分别为棱 \(AB\)、\(BC\) 的中点，\(M\) 为棱 \(AA_{1}\) 上的点，二面角 \(M - DE - A\) 为 \(30^{\circ}\).

\bitmapfigure[max width=0.38\linewidth,max totalheight=4.8cm]{img/2007/liaoning_science/q18_fig1.png}

(1) 证明： \(A_{1}B_{1} \perp C_{1}D\);

(2) 求 \(MA\) 的长，并求点 \(C\) 到平面 \(MDE\) 的距离.
\end{problem}

\begin{problem}
已知函数 \(f(x) = \sin \left(\omega x + \frac{\pi}{6}\right) + \sin \left(\omega x - \frac{\pi}{6}\right) - 2 \cos^2 \frac{\omega x}{2}, x \in \R\)（其中 \(\omega > 0\)）.

(1) 求函数 \(f(x)\) 的值域;

(2) 若函数 \(y = f(x)\) 的图象与直线 \(y = -1\) 的两个相邻交点的距离为 \(\frac{\pi}{2}\)，求函数 \(y = f(x)\) 的单调区间.
\end{problem}

\begin{problem}
已知正三角形 \(OAB\) 的三个顶点都在抛物线 \(y^{2} = 2x\) 上，其中 \(O\) 为坐标原点，设圆 \(C\) 是 \(\triangle OAB\) 的外接圆（点 \(C\) 为圆心）.

(1) 求圆 \(C\) 的方程;

(2) 设圆 \(M\) 的方程为 \((x - 4 - 7\cos \theta)^2 + (y - 7\sin \theta)^2 = 1\)，过圆 \(M\) 上任意一点 \(P\) 分别作圆 \(C\) 的两条切线 \(PE\)、\(PF\)，切点为 \(E\)、\(F\)，求 \(\overrightarrow{CE} \cdot \overrightarrow{CF}\) 的最大值和最小值.
\end{problem}

\begin{problem}
已知数列 \(\{a_{n}\}\)、\(\{b_{n}\}\) 与函数 \(f(x)\)、\(g(x)\)，\(x \in \R\) 满足条件：\(b_{1} = b\)，\(a_{n} = f(b_{n}) = g(b_{n+1})\)（\(n \in \mathbf{N}^{*}\)）.

(1) 若 \(f(x) = tx + 1\)（\(t \neq 0, t \neq 2\)），\(g(x) = 2x\)，\(f(b) \neq g(b)\)，且 \(\lim_{n \to \infty} a_n\) 存在，求 \(t\) 的取值范围，并求 \(\lim_{n \to \infty} a_n\)（用 \(t\) 表示）;

(2) 若函数 \(y = f(x)\) 为 \(\R\) 上的增函数，\(g(x) = f^{-1}(x)\)，\(b = 1\)，\(f(1) < 1\)，证明对任意的 \(n \in \mathbf{N}^{*}\)，\(a_{n+1} < a_n\).
\end{problem}

\begin{problem}
已知函数 \(f(x) = e^{2x} - 2t(e^{x} + x) + x^{2} + 2t^{2} + 1\)，\(g(x) = \frac{1}{2} f'(x)\).

(1) 证明：当 \(t < 2\sqrt{2}\) 时，\(g(x)\) 在 \(\R\) 上是增函数;

(2) 对于给定的闭区间 \([a, b]\)，试说明存在实数 \(k\)，当 \(t > k\) 时，\(g(x)\) 在闭区间 \([a, b]\) 上是减函数;

(3) 证明：\(f(x) \geqslant \frac{3}{2}\).
\end{problem}
